\documentclass{css}
%\documentclass[english]{css}

\usepackage[dvipdfmx]{graphicx}
\usepackage{latexsym}
\usepackage{amsmath}

\def\|{\verb|}

\newcommand{\cssyear}[0]{2026}
\newcommand{\cssname}[0]{CSS 2026}

% ============================================================
% NOTE (著者向け・提出時は削除):
% 本稿の全数値は census2/（go.work・toolchain の偽EMPTY_GT を修正した
% 公正な単一クローン計測）に基づく。旧 census/（v1）とは評価対象数・各表が
% 微差で異なる（定性結論＝GT定義による優劣逆転は両者で不変）。
% 出典: census2/SUMMARY_census.md, census2/manifest.csv, census2/metrics.csv
% ============================================================

\begin{document}

\title{Go言語を対象とした複数のGround Truth定義による\\
SBOMツール精度評価}
\etitle{Accuracy Evaluation of SBOM Generation Tools for Go\\
under Multiple Ground Truth Definitions}

\affiliate{Lab}{○○大学 情報科学研究室\\
○○ University}

\author{安岡 瑞希}{Mizuki Yasuoka}{Lab}[mizuki@example.ac.jp]
%\author{中村 修}{Osamu Nakamura}{Lab}

\begin{abstract}
本稿では，Go言語プロジェクトを対象に，SBOM生成ツールの精度を複数の
Ground Truth（GT）定義のもとで評価する．awesome-goから取得した
リポジトリのうち評価可能な1{,}528個を対象に，4つのツール（Syft，Trivy，
cdxgen，cyclonedx-gomod）を評価し，GTをimported / imported+test / all
の3定義で比較した結果，どのGT定義で測るかによってツールの優劣が逆転する
ことを示した．さらに，この逆転は各ツールが依存を読み取る源とGT定義の
一致度で説明でき，本評価はツールの絶対的な優劣ではなくGT選択への感度を
示すものであることを論じる．
\end{abstract}

\begin{jkeyword}
SBOM，ソフトウェアサプライチェーン，Go言語，精度評価，Ground Truth
\end{jkeyword}

\begin{eabstract}
In this paper, we evaluate the accuracy of SBOM generation tools for
Go projects under multiple ground-truth (GT) definitions. Evaluating
four tools (Syft, Trivy, cdxgen, and cyclonedx-gomod) on 1{,}528
evaluable Go projects from awesome-go under three GT definitions
(imported / imported+test / all), we show that the ranking of the
tools reverses depending on which GT definition is used. We further
argue that this reversal is explained by how closely each tool's
information source matches each GT definition, so the study measures
sensitivity to the choice of GT rather than absolute tool quality.
\end{eabstract}

\begin{ekeyword}
SBOM, software supply chain, Go language, accuracy evaluation, ground truth
\end{ekeyword}

\maketitle

%% ============================================================
\section{はじめに}
%% ============================================================
\begin{itemize}
\item ソフトウェアサプライチェーン攻撃が急増しており，対策として，
ソフトウェアの部品一覧であるSBOMが注目されている．
\item SBOMの利点は一覧が正確なときにのみ得られる．不正確だと，危険な
依存の見落とし，または未使用の依存の誤報を招く．
\item ツールの正確性は，正解一覧であるGround Truth（GT）と突き合わせて
評価する．しかしGTの定義次第で正解の範囲は変わり，評価結果も変わりうる．
\item PythonやJava等ではGTベース評価が行われているが，いずれも単一の
GT定義に基づく．とくにGo言語ではGTベース評価自体が行われていない．
\item 本稿の貢献：(1) Goを対象にGTベース精度評価を初めて実施
（1{,}528プロジェクト，4ツール），(2) GTを3定義で比較し定義により
ツール優劣が逆転することを定量的に提示，(3) 逆転の要因を各ツールが依存を
読み取る源（宣言か，コンパイルグラフか）の違いから説明．
\end{itemize}

%% ============================================================
\section{背景}
%% ============================================================

\subsection{SBOMとその正確性}
\begin{itemize}
\item SBOMは，ソフトが使う部品（依存ライブラリ）の一覧表である．
\item 脆弱性対応やライセンス調査に用いる．一覧が不正確だと実務判断を誤る．
\item 正確性の測定には，答え合わせ用の正解一覧（GT）が必要である．
\end{itemize}

\subsection{Goの依存の単位}
\begin{itemize}
\item Goの依存は、モジュールがパッケージをまとめた階層構造をなす。SBOMでは、パッケージではなく、そのパッケージが属するモジュールを依存の単位として扱う。
\item コードは他のパッケージをimportして利用する。あるプロジェクトが必要とする依存は、このimport関係を辿ることで特定できる。
\item 標準ライブラリは依存に数えず、外部から取得するパッケージが属するモジュールをSBOMの対象とする。
\item テストコードでのみ使われる依存（テスト依存）は、本番のコードでは使われない。この区別は後の正解データの定義に関わる。
\end{itemize}

\subsection{go.mod と go.sum}
\begin{itemize}
\item go.modは，どの依存のどのバージョンを使うかをモジュール単位で
記述する．直接依存と間接依存（\|// indirect|）を含み，実際には使わない
依存が含まれることもある．
\item go.sumは各依存（間接含む）のハッシュを記録する．過去に関わった
依存・バージョンも残るため，go.modより広く，今使わないものも含む．
\item go.modは，\|go mod init|で初期化し，\|go mod tidy|を実行すると，
依存先が要求する間接依存も一括で自動追加される（実コードでの使用有無は
問わない）．このため「使わない依存」が生じる．
\end{itemize}

\begin{figure}[tb]
% \includegraphics[width=\columnwidth]{fig_gomod_flow.pdf}
\caption{go.modの生成過程（ginの例）}
\ecaption{How go.mod is generated (gin example).}
\label{fig:gomod}
\end{figure}

%% ============================================================
\section{Ground Truthの3つの定義}
%% ============================================================
\begin{itemize}
\item 正解データGTを\tabref{tab:gtdef}の3通りで定義する．いずれも
\|go list|により，rootから\|GOOS=linux|で生成する．
\item GT-importedは，非testの実コンパイル依存だけを集めたものである．
\item GT-imported+test（impT）は，それにtest専用依存を加えたものである．
\item GT-allは，build list全体（直接・間接・未使用を含む）であり，最も
広い定義である．
\item 3定義は\figref{fig:venn}のように，概ね imported $\subseteq$ impT
$\subseteq$ all の包含関係にある．
\end{itemize}

\begin{table}[tb]
\caption{正解データ（GT）の3つの定義}
\ecaption{Three definitions of the ground truth.}
\label{tab:gtdef}
\hbox to\hsize{\hfil
\begin{tabular}{ll}\hline\hline
GT定義 & 生成コマンド（要点） \\\hline
GT-imported & \|go list -deps| \\
GT-imported+test & \|go list -deps -test| \\
GT-all & \|go list -m all| \\\hline
\end{tabular}\hfil}
\end{table}

\begin{figure}[tb]
% \includegraphics[width=\columnwidth]{fig_venn.pdf}
\caption{3つのGT定義の包含関係}
\ecaption{Inclusion relation among the three GT definitions.}
\label{fig:venn}
\end{figure}

%% ============================================================
\section{関連研究}
%% ============================================================
\begin{itemize}
\item GTベースのSBOM精度評価は，複数の言語で行われてきた
（\tabref{tab:related}）．
\item いずれの研究も単一のGT定義に基づいており，GT定義の選択が評価に
与える影響は検討されていない．
\item 本研究は，Goを対象にGTベース評価を初めて行い，かつGT定義を3通りに
変えて結果の変化を比較する点に新規性がある．
\end{itemize}

\begin{table*}[tb]
\caption{各言語における既存のGTベース評価と本研究の位置づけ}
\ecaption{Existing GT-based evaluations per language and our positioning.}
\label{tab:related}
\hbox to\hsize{\hfil
\begin{tabular}{lll}\hline\hline
言語 & 研究 & 正解（GT）の作り方 \\\hline
Python & Yu et al. (2024) & 実際にインストールして入ったものを正解 \\
Python & Reality Check (2025) & ロックファイルを正解 \\
Java & Balliu et al. (2023) & Maven依存ツリー出力で一覧 \\
JavaScript & Rabbi et al. (2024) & \|npm list| で一覧 \\
Rust & Reality Check (2025) & ロックファイルを正解 \\
Ruby & Reality Check (2025) & ロックファイルを正解 \\
PHP & Reality Check (2025) & ロックファイルを正解 \\
C / C++ & TPL Dependency研究 & 人手でコードを読み正解作成 \\
Go & 本研究 & \|go list| でimport依存を集め正解 \\\hline
\end{tabular}\hfil}
\end{table*}

%% ============================================================
\section{評価手法}
%% ============================================================

\subsection{研究目的とリサーチクエスチョン}
\begin{itemize}
\item 主目的は，Goプロジェクトを対象にGTベースで4ツールの精度を評価する
ことである．
\item RQ1：GT定義（all / imported / imported+test）を変えると，ツールの
優劣は変わるか．
\item RQ2：変わるとすれば，その原因は何か．
\end{itemize}

\subsection{対象データセット}
\begin{itemize}
\item awesome-goから2{,}723件を取得・記録した．
\item statusの内訳は，CLONE\_FAIL（取得失敗）13件，EMPTY\_GT（外部依存を
持たず正解GTが空）1{,}182件，OK（imported依存が1件以上あり評価対象）
1{,}528件である．
\item 本研究の評価対象は，このOKの1{,}528件である．
\end{itemize}

\subsection{評価対象ツールと実行環境}
\begin{itemize}
\item Syft v1.46.0，Trivy v0.72.0，cdxgen 12.7.1（Node.js v22.22.2），
cyclonedx-gomod v1.10.0を用いた．
\item GTの生成には go1.26.5（\|GOTOOLCHAIN=local|，\|GOFLAGS=-mod=mod|）
を用い，\|GOOS=linux|で実行した．GTもツールも同一のクローン上で同時に
生成・照合した（全リポジトリ1回ずつのクローン）．
\item 依存はモジュールパス単位で照合する（name一致）．本編ではname一致を
主とし，version一致は付録で補足する．
\item GT生成の環境設定（\|-mod=mod|の自動解除，新Go要求リポジトリの
toolchain取得，\|-e|の扱い）は付録\ref{app:gtimpl}に詳述する．これらは
生成コマンドの式は変えず，同じ式がより多くのリポジトリで成功するように
する設定である．
\end{itemize}

\subsection{各ツールが依存を読み取る源}
\begin{itemize}
\item SyftとTrivyは，go.mod・go.sum・ツリー内の別go.modを静的に読む
（ビルド不要）．go.sumは最も広い集合のため過剰報告が多く，precisionが
低くrecallが高い傾向をもつ．
\item cdxgenは\|go list -deps|と\|go mod graph|を用い，実コンパイル
グラフに最も近い依存を列挙する．
\item cyclonedx-gomodは\|go list -m all|を\|go mod why|で到達可能性
フィルタし，到達可能なモジュールだけを残す公式ツールである．
\item 統一的には，Syft/Trivyは「宣言（go.sum）」を，cdxgen/
cyclonedx-gomodは「コンパイルグラフ・到達可能性」を報告する設計の違い
であり，バグではない．
\end{itemize}

\begin{table}[tb]
\caption{各SBOMツールが読み取る源と近いGT}
\ecaption{Source read by each tool and the closest GT.}
\label{tab:tools}
\hbox to\hsize{\hfil
\begin{tabular}{lll}\hline\hline
ツール & 源 & 近いGT \\\hline
Syft & go.mod ＋ go.sum & GT-all \\
Trivy & go.mod ＋ go.sum & GT-all \\
cdxgen & \|go list -deps| & GT-imported \\
cyclonedx-gomod & \|go list -m all| ＋ 到達性 & GT-imported \\\hline
\end{tabular}\hfil}
\end{table}

\subsection{評価指標}
\begin{itemize}
\item ツール出力とGTの対応（TP/FP/FN/TN）から，precision，recall，F1を
算出する（式(\ref{eq:prec})〜(\ref{eq:f1})）．
\item 各GT定義に対し，リポジトリごとにP/R/F1を算出し，評価可能な
リポジトリで平均する（macro平均）．
\item 空GTのリポジトリはrecallが定義できないため除外し，ツールが出力
できなかったリポジトリはNAとして扱う．
\end{itemize}

\begin{equation}
\mathrm{precision} = \frac{\mathrm{TP}}{\mathrm{TP}+\mathrm{FP}}
\label{eq:prec}
\end{equation}
\begin{equation}
\mathrm{recall} = \frac{\mathrm{TP}}{\mathrm{TP}+\mathrm{FN}}
\label{eq:rec}
\end{equation}
\begin{equation}
\mathrm{F1} = 2 \times
\frac{\mathrm{precision}\times\mathrm{recall}}
{\mathrm{precision}+\mathrm{recall}}
\label{eq:f1}
\end{equation}

%% ============================================================
\section{結果}
%% ============================================================

\subsection{GT定義によるツール優劣の逆転}
\begin{itemize}
\item 各ツールのmacro平均P/R/F1（name一致）を3つのGT定義ごとに
\tabref{tab:result}に示す．
\item GT-allでは，SyftとTrivyのF1が最も高く（73.0，72.5），cdxgenと
cyclonedx-gomodは低い（52.8，56.5）．
\item GT-importedでは順位が入れ替わり，cdxgenとcyclonedx-gomodの
F1が最も高くなる（93.2，94.5）．
\item すなわち，どのGT定義で測るかによってツールの優劣が逆転する．
\end{itemize}

\begin{table}[tb]
\caption{各ツールの精度（macro平均・name一致，all / imported /
imported+test，\%）}
\ecaption{Accuracy per tool (macro, name-match; all/imp/impT, \%).}
\label{tab:result}
\hbox to\hsize{\hfil
\begin{tabular}{lll}\hline\hline
ツール & 指標 & 値 \\\hline
Syft & prec. & 92.6 / 55.6 / 68.6 \\
     & rec.  & 67.1 / 99.9 / 99.9 \\
     & F1    & 73.0 / 67.0 / 78.3 \\\hline
Trivy & prec. & 92.3 / 59.1 / 72.9 \\
      & rec.  & 67.0 / 99.9 / 99.9 \\
      & F1    & 72.5 / 68.4 / 79.7 \\\hline
cdxgen & prec. & 94.5 / 91.4 / 92.2 \\
       & rec.  & 42.4 / 99.8 / 82.4 \\
       & F1    & 52.8 / 93.2 / 82.3 \\\hline
cyclonedx & prec. & 99.9 / 92.4 / 92.6 \\
-gomod    & rec.  & 43.5 / 99.1 / 81.0 \\
          & F1    & 56.5 / 94.5 / 82.5 \\\hline
\end{tabular}\hfil}
\end{table}

\subsection{テスト依存の影響}
\begin{itemize}
\item GT-imported+testでは，テスト依存を正解に含めることで
SyftとTrivyのF1が上昇し（67.0→78.3，68.4→79.7），cdxgenと
cyclonedx-gomodは低下する（93.2→82.3，94.5→82.5）．
\item テスト依存をGTに含めるか否かも，ツールの評価結果に影響する．
\end{itemize}

\subsection{母数差に対するロバストネス}
\begin{itemize}
\item ツールごとに出力失敗（NA）の件数が異なり，有効評価repo数は
Syft 1{,}519，Trivy 1{,}509，cdxgen 1{,}493，cyclonedx-gomod 1{,}501と
差がある（付録\ref{app:raw}）．
\item この母数差が順位を歪めていないことを確認するため，4ツールすべてが
成功した共通repo集合（$n=1{,}467$）に限定して再計算した（\tabref{tab:robust}）．
\item 値の変化はいずれも1ポイント未満であり，順位（cyclonedx-gomod
$>$ cdxgen $\gg$ Trivy $\approx$ Syft）は完全に一致した．
\item よって，ツール間の母数差は本研究の結論に影響しない．
\end{itemize}

\begin{table}[tb]
\caption{母数差のロバストネス検証（imported，name一致，macro-F1，\%）}
\ecaption{Robustness to differing denominators
(imported, name-match, macro-F1, \%).}
\label{tab:robust}
\hbox to\hsize{\hfil
\begin{tabular}{lrr}\hline\hline
ツール & 各自の有効集合 & 共通集合（$n=1{,}467$） \\\hline
Syft & 67.0 & 67.4 \\
Trivy & 68.4 & 68.8 \\
cdxgen & 93.2 & 93.3 \\
cyclonedx-gomod & 94.5 & 94.7 \\\hline
\end{tabular}\hfil}
\end{table}

%% ============================================================
\section{考察}
%% ============================================================

\subsection{逆転が生じる理由}
\begin{itemize}
\item ツールが読み取る源（依存の範囲）と正解GTの範囲がずれた分だけ，
スコアは落ちる．
\item GT-allは build list（\|go list -m all|）に近く，go.sumを広く読む
Syft / Trivyが近くなりF1が高い．
\item GT-importedは実importに近く，コンパイルグラフを読むcdxgen /
cyclonedx-gomodが近くなりF1が高い．
\item よって，同じツール群でも，どのGTで測るかで優劣が逆転する．これは
ツールのバグではなく，各ツールが「宣言」を報告するか「到達可能性」を
報告するかという設計選択の違いに由来する．
\end{itemize}

\subsection{評価の枠組みとしての含意}
\begin{itemize}
\item GT-importedの生成は\|go list -deps|であり，これはcdxgenが内部で
用いる機構とほぼ同一である．同様にGT-allはSyft/Trivyが読むgo.sumに近い．
したがって「各ツールが自分の源に一致するGTで高得点となる」ことは，
半ば構成上の必然である．
\item この意味で，本評価が示すのはツールの絶対的な優劣ではなく，
\textbf{GT定義の選択に対する評価結果の感度}である．どのツールが「正しい」
かはGTをどう定義するかに依存し，中立なGTは存在しない．
\end{itemize}

\subsection{実務への含意}
\begin{itemize}
\item 用途に応じてGTを選ぶ必要がある．脆弱性管理では実際に使う依存を
見たいのでimported，ライセンス調査では関わる依存全部を見たいのでallが
適する．
\item 用途に合わないGTで測ると，ツール選定を誤るおそれがある．
\end{itemize}

\subsection{限界}
\begin{itemize}
\item 対象はawesome-goに収録された公開リポジトリに限られ，企業の非公開
プロジェクト等への一般化には注意を要する．
\item 正解GTを\|go list|に基づいて定義しており，このGTの作り方自体が
一つの立場である．とくにcdxgen/cyclonedx-gomodは内部で\|go list|を用いる
ため，これらのGTに対する高い一致は独立な精度というより整合性の確認に
近い側面をもつ．別のGT定義を採用すれば結果が変わる可能性がある．
\item 照合はname一致（モジュールパス単位）を主としており，バージョンの
厳密な一致は付録での補足にとどまる．
\item 各ツールは既定に近い設定で1回実行しており，オプションやモードの
違いによる感度は検討していない．
\end{itemize}

%% ============================================================
\section{おわりに}
%% ============================================================
\begin{itemize}
\item Goを対象に，GTベースでSBOMツールの精度を大規模に評価した
（1{,}528プロジェクト，4ツール）．
\item GTを3定義で比較し，定義によってツールの優劣が変わることを定量的に
示すとともに，その要因を各ツールが依存を読み取る源の違いから説明した．
\item 今後の課題として，偽陽性の要因分析の深掘り，対象言語・データセット
の拡張，およびツール設定の感度分析が挙げられる．
\end{itemize}

%\begin{acknowledgment}
%本研究の遂行にあたりご指導いただいた中村 修 教授に感謝する．
%\end{acknowledgment}

\begin{thebibliography}{9}
\bibitem{yu2024}
S. Yu, et al.:
On the Correctness of Metadata-Based SBOM Generation:
A Differential Analysis Approach,
Proc. DSN, pp.29--36 (2024).

\bibitem{balliu2023}
M. Balliu, et al.:
Challenges of Producing Software Bill of Materials for Java,
IEEE Security \& Privacy (2023).

\bibitem{rabbi2024}
M. F. Rabbi, et al.:
SBOM Generation Tools Under Microscope:
A Focus on the npm Ecosystem,
Proc. SAC, pp.1233--1241 (2024).

\bibitem{zhou2026}
L. Zhou, et al.:
A Reality Check on SBOM-based Vulnerability Management,
Proc. CODASPY (2026).

\bibitem{tang2022}
W. Tang, et al.:
Towards Understanding Third-party Library Dependency in
C/C++ Ecosystem,
Proc. ASE (2022).

\bibitem{golist}
The Go Authors:
Command go (go list),
\|https://go.dev/cmd/go|.

\bibitem{gomod}
The Go Authors:
Go Modules Reference,
\|https://go.dev/ref/mod|.

\bibitem{syft}
Anchore: Syft, \|https://github.com/anchore/syft|.

\bibitem{trivy}
Aqua Security: Trivy, \|https://github.com/aquasecurity/trivy|.

\bibitem{cdxgen}
OWASP CycloneDX: cdxgen,
\|https://github.com/CycloneDX/cdxgen|.

\bibitem{gomodtool}
OWASP CycloneDX: cyclonedx-gomod,
\|https://github.com/CycloneDX/cyclonedx-gomod|.
\end{thebibliography}

%% ============================================================
%% 付録（本文と合わせて全体8ページ以内）
%% ============================================================
\appendix

%% ------------------------------------------------------------
\section{version一致・micro集計の精度}
%% ------------------------------------------------------------
本文\tabref{tab:result}はmacro平均・name一致である．参考として，
version一致（\tabref{tab:macro-ver}）と，プール合計から算出したmicro
集計（\tabref{tab:micro-name}）を示す．いずれも優劣逆転の傾向は保存
される．

\begin{table}[tb]
\caption{macro平均・version一致（all / imp / impT，\%）}
\label{tab:macro-ver}
\hbox to\hsize{\hfil
\begin{tabular}{lll}\hline\hline
ツール & 指標 & 値 \\\hline
Syft & prec. & 91.1 / 54.8 / 67.8 \\
     & rec.  & 66.6 / 99.9 / 99.9 \\
     & F1    & 72.1 / 66.3 / 77.5 \\\hline
Trivy & prec. & 90.9 / 58.5 / 72.1 \\
      & rec.  & 66.4 / 99.9 / 99.8 \\
      & F1    & 71.6 / 67.8 / 79.1 \\\hline
cdxgen & prec. & 93.5 / 90.8 / 91.6 \\
       & rec.  & 41.9 / 99.7 / 82.3 \\
       & F1    & 52.2 / 92.7 / 81.8 \\\hline
cyclonedx & prec. & 99.9 / 92.3 / 92.6 \\
-gomod    & rec.  & 43.5 / 99.1 / 81.0 \\
          & F1    & 56.5 / 94.5 / 82.4 \\\hline
\end{tabular}\hfil}
\end{table}

\begin{table}[tb]
\caption{micro集計・name一致（all / imp / impT，\%）}
\label{tab:micro-name}
\hbox to\hsize{\hfil
\begin{tabular}{lll}\hline\hline
ツール & 指標 & 値 \\\hline
Syft & prec. & 84.9 / 53.0 / 57.5 \\
     & rec.  & 50.0 / 99.8 / 99.8 \\
     & F1    & 62.9 / 69.2 / 73.0 \\\hline
Trivy & prec. & 84.4 / 54.4 / 59.2 \\
      & rec.  & 48.1 / 99.8 / 99.8 \\
      & F1    & 61.3 / 70.4 / 74.3 \\\hline
cdxgen & prec. & 88.2 / 82.7 / 83.6 \\
       & rec.  & 33.1 / 99.8 / 92.5 \\
       & F1    & 48.2 / 90.5 / 87.8 \\\hline
cyclonedx & prec. & 99.8 / 85.6 / 85.7 \\
-gomod    & rec.  & 36.0 / 98.9 / 91.1 \\
          & F1    & 52.9 / 91.8 / 88.3 \\\hline
\end{tabular}\hfil}
\end{table}

%% ------------------------------------------------------------
\section{TP / FP / FN プール合計（name一致）}
\label{app:raw}
%% ------------------------------------------------------------
\begin{table}[tb]
\caption{TP / FP / FN プール合計（name一致）}
\label{tab:raw}
\hbox to\hsize{\hfil
\begin{tabular}{llll}\hline\hline
ツール & all & imported & impT \\
 & (TP/FP/FN) & (TP/FP/FN) & (TP/FP/FN) \\\hline
Syft & 76{,}929 / 13{,}673 / 77{,}030
     & 47{,}981 / 42{,}621 / 120
     & 52{,}136 / 38{,}466 / 123 \\
Trivy & 71{,}376 / 13{,}198 / 77{,}001
      & 45{,}970 / 38{,}604 / 109
      & 50{,}075 / 34{,}499 / 112 \\
cdxgen & 47{,}975 / 6{,}413 / 96{,}881
       & 45{,}006 / 9{,}382 / 77
       & 45{,}442 / 8{,}946 / 3{,}705 \\
cyclonedx & 54{,}433 / 130 / 96{,}653
-gomod    & 46{,}700 / 7{,}863 / 526
          & 46{,}776 / 7{,}787 / 4{,}578 \\\hline
\end{tabular}\hfil}
\end{table}

\begin{table}[tb]
\caption{有効評価repo数とNA数（評価対象1{,}528のうち）}
\label{tab:valid}
\hbox to\hsize{\hfil
\begin{tabular}{lrr}\hline\hline
ツール & 有効 & NA \\\hline
Syft & 1{,}519 & 9 \\
Trivy & 1{,}509 & 19 \\
cdxgen & 1{,}493 & 35 \\
cyclonedx-gomod & 1{,}501 & 27 \\\hline
\end{tabular}\hfil}
\end{table}

%% ------------------------------------------------------------
\section{偽陽性の要因分析（imported基準・name）}
%% ------------------------------------------------------------
\begin{itemize}
\item importedを基準としたときの偽陽性（FP）の要因内訳を
\tabref{tab:fp}に示す．
\item Syft / Trivyでは残余（go.sum残骸＋兄弟go.mod）が最大の要因であり，
go.sumを広く読むことに起因する．
\item cyclonedx-gomodではindirect未使用が最大の要因であり，到達可能性
フィルタを経てなお残る間接依存に起因する．
\end{itemize}

\begin{table}[tb]
\caption{FP要因バケツ（imported基準・name，\%）}
\label{tab:fp}
\hbox to\hsize{\hfil
\begin{tabular}{lrrrrr}\hline\hline
ツール & test & 他OS & direct未使用 & indirect未使用 & 残余 \\\hline
Syft & 9.7 & 2.1 & 0.8 & 16.4 & 71.0 \\
Trivy & 10.6 & 2.3 & 0.8 & 17.0 & 69.2 \\
cdxgen & 4.6 & 1.6 & 0.9 & 9.6 & 83.3 \\
cyclonedx-gomod & 1.0 & 11.5 & 3.1 & 62.6 & 21.8 \\\hline
\end{tabular}\hfil}
\end{table}

%% ------------------------------------------------------------
\section{GT生成の実装上の注意}
\label{app:gtimpl}
%% ------------------------------------------------------------
GTを全リポジトリで確実に生成するため，以下の環境設定を行った．いずれも
\|go list|の\emph{式}（依存の集め方）は変えず，同じ式がより多くの
リポジトリで成功するようにするものである．誤って正解GTが空と判定される
（偽EMPTY\_GT）ことによる取りこぼしを防ぐ．
\begin{itemize}
\item \textbf{\texttt{-e}フラグ：} \|go list|に\texttt{-e}を付けると，
壊れたパッケージがあってもエラーで止まらず，解決できる依存の列挙を続ける．
正常なモジュールでは\texttt{-e}の有無で出力（依存一覧）は同一であり，
異なるのは終了コードのみである．\texttt{-e}が結果に影響するのは，
\texttt{-e}なしだと列挙が全て落ちて空になる稀なケースに限られ，そこでは
取りこぼしを減らす方向に働く（偽の依存を追加しない）．
\item \textbf{vendor対応（\texttt{-mod=mod}）：} \texttt{vendor/}
ディレクトリを持つリポジトリは，既定の\texttt{-mod=vendor}では
\|go list -m all|が失敗し，正解GTが空と誤判定される．\texttt{-mod=mod}を
付けることでこれを回避し，vendored repoの取りこぼしを防ぐ．
\item \textbf{ワークスペース対応（\texttt{go.work}）：}
\texttt{go.work}を持つリポジトリはworkspace modeとなり，複数モジュールを
まとめて扱うため，goは複数のgo.modの自動編集を許さない．このため
workspace modeでは\texttt{-mod}に\texttt{readonly}か\texttt{vendor}しか
指定できず，前項の\texttt{-mod=mod}はエラーとなって\|go list|が全て失敗し，
偽EMPTY\_GTを生じる．そこで\texttt{go.work}を検出したリポジトリでは
\texttt{-mod=mod}を自動的に解除し，workspace既定（readonly）で解析する
（例：etcd，kubernetes系，pomerium等がこれに該当）．
\item \textbf{新しいGoを要求するリポジトリ：} 基準は
go1.26.5＋\texttt{GOTOOLCHAIN=local}（per-repoのtoolchainを取得しない
設定）である．go.modが基準より新しいGoを要求する場合（例：\texttt{go
1.27rc2}）はビルド不可で\|go list|が失敗するため，該当リポジトリのみ
\texttt{GOTOOLCHAIN=auto}として必要なtoolchainを取得して計測した．
\item \textbf{正しく空となる場合：} 外部依存を持たず標準ライブラリのみを
用いるリポジトリは，GT-importedが正しく空となりEMPTY\_GT（評価対象外）と
なる．自モジュール自身は依存に数えない．
\end{itemize}

\end{document}
